\section{Discussion}
Our results support assigning explicit strategic decisions to a separate controller when the negotiation rules can be specified. The planner uses beliefs, action costs, and future consequences to improve decisions under the same language conditions. Distillation transfers part of this behavior to a policy that acts without search or further language model calls. Its improvement over base PPO shows the value of planner supervision in this setting. Its remaining gap to Sol medium and small advantage over the heuristic distinguish exact planning from learned control.

The interventions suggest that asking more questions is not the main source of the learned policy's improvement. Signals change acceptance incentives, while responses to offers can update beliefs without a separate question. Useful information therefore depends on what the agent can subsequently do with it and what obtaining it costs. Discourse analysis adds a complementary requirement: the chosen action must be expressed with the intended conditions and commitment strength. A conditional threat to withdraw and an actual withdrawal can lead to different transitions. Strategic evaluation should therefore connect utility, information use, and the meaning of the resulting utterance.

\section{Conclusion}
We studied strategic control for a lower-cost language model in a diplomatic negotiation simulation grounded in real documents. With utterances held fixed, belief-state planning achieved the highest utility in the matched comparison, and planner-distilled PPO improved on base PPO without further model calls for action selection. Interventions linked the learned policy's gains mainly to signals and updated beliefs. Dialogue traces further showed how information use and commitment wording shape negotiation paths. These findings identify explicit strategic control as a useful component of negotiation systems and show how controlled comparisons can explain its benefits.
